\documentclass[11pt,letterpaper]{article}
\usepackage[left=.95in,right=.95in,top=.82in,bottom=.85in]{geometry}
\usepackage[T1]{fontenc}
\usepackage{lmodern}
\usepackage{amsmath,amssymb,amsthm,mathtools}
\usepackage[expansion=false]{microtype}
\usepackage{xcolor,booktabs}
\definecolor{ink}{HTML}{17485D}
\usepackage[colorlinks=true,linkcolor=ink,citecolor=ink,urlcolor=ink]{hyperref}
\usepackage{fancyhdr}
\pagestyle{fancy}\fancyhf{}\fancyhead[L]{\small\itshape Fourier measurements of Dirichlet zeros}\fancyhead[R]{\small Research note}\fancyfoot[C]{\thepage}
\setlength{\headheight}{14pt}
\newtheorem{theorem}{Theorem}
\newtheorem{proposition}[theorem]{Proposition}
\newtheorem{corollary}[theorem]{Corollary}
\newcommand{\F}{\mathbb F}
\newcommand{\Z}{\mathbb Z}
\newcommand{\Q}{\mathbb Q}
\newcommand{\C}{\mathbb C}
\newcommand{\R}{\mathbb R}
\newcommand{\eps}{\varepsilon}
\newcommand{\cB}{\mathcal B}
\newcommand{\muE}{\boldsymbol\mu_E}
\newcommand{\abs}[1]{\lvert #1\rvert}
\DeclareMathOperator{\ord}{ord}
\setlength{\emergencystretch}{1.8em}
\raggedbottom
\hypersetup{pdftitle={What zero differences forget about Dirichlet characters},pdfauthor={OpenAI ChatGPT}}
\begin{document}
\thispagestyle{plain}
\noindent{\small\color{ink}\sffamily ANALYTIC NUMBER THEORY\hfill 26 NOVEMBER 2026}
\vspace{9pt}
\begin{center}
{\LARGE\bfseries What zero differences forget\\[3pt] about Dirichlet characters}\par
\end{center}
\begin{abstract}
We classify the balanced Fourier measurements of the pooled positive zero ordinates of a product of primitive Dirichlet $L$-functions. They determine, at each unramified prime, the multiset of twist values up to a rotation that can depend on the prime. For quadratic twists this is exactly equality of the multisets of pairwise character quotients. At conductor $32045$, one measurement class contains exactly 28 common-twist classes of six-factor products in an explicit family. Two of them differ by at least $(\sqrt2\log 2/\pi)T-O(\log T)$ zero occurrences up to height $T$. Despite this discrepancy, their balanced measurements agree asymptotically even when the frequencies grow to $(2-\delta)\log T$. For twists of odd prime exponent, a single complex cubic measurement per residue class recovers its local multiset up to rotation. When the packet size is coprime to the exponent, these local rotations come from one global twist.
\end{abstract}

\section{The measurement}
A \emph{packet} is a finite nonempty multiset of characters. Pool the nontrivial zeros of its $L$-functions, adding multiplicities, and write $Z_A(T)$ for the resulting multiset of ordinates in $(0,T]$. An observer who knows only differences can measure, for $d\ge2$ and real $t_j\ne0$ with $\sum_j t_j=0$,
\begin{align}
 S_A(t;T)&=\frac1T\sum_{\gamma\in Z_A(T)}e^{it\gamma},
 &\cB_A(\mathbf t;T)&=\prod_{j=1}^d S_A(t_j;T). \label{eq:observable}
\end{align}
Indeed, expansion gives a sum over all ordered tuples. Each phase is
$\exp(i\sum_{j<d}t_j(\gamma_j-\gamma_d))$. A common translation of the observed heights cancels exactly. The normalization by $T$ retains the contribution of prime powers to a nonzero Fourier mode. Division by the number of zeros would wash it out, since that number has order $T\log T$.

For $d=2$ the measurement is $|S_A(t;T)|^2$. One might expect triples and longer tuples to recover phases lost upon squaring. For quadratic twists, even all orders leave the same ambiguity. The example in Section~\ref{sec:six} has equal limits for every such tuple, while its discrepancies occur in every sufficiently long logarithmic height interval. Section~\ref{sec:uniform} makes the equality uniform over a growing range of frequencies. These results concern the normalization and Fourier tests in \eqref{eq:observable}. Local zero spacings and statistics normalized on the mean-spacing scale pose different questions.

\paragraph{The two packets at a glance.}
The 16 characters in the example have four-bit labels. Rows give the first two bits, columns the last two, and each 1 selects a factor of the product. Each table selects six factors. Their pairwise character quotients have exactly the same multiplicities.
\[
\begin{array}{c|cccc}
 A &00&01&10&11\\\hline
 00&0&0&0&1\\01&0&0&0&1\\10&0&0&0&1\\11&1&1&1&0
\end{array}
\qquad\qquad
\begin{array}{c|cccc}
 B &00&01&10&11\\\hline
 00&0&0&0&1\\01&0&0&0&1\\10&0&0&1&0\\11&1&1&0&1
\end{array}
\]

\newpage
\section{From ordinates to roots of unity}\label{sec:prime}
Fix a modulus $Q$, its unit group $G=(\Z/Q\Z)^\times$, a character $\psi$ modulo $Q$, and a subgroup $H\le\widehat G$ of exponent $E$. A multiplicity function $A:H\to\Z_{\ge0}$ denotes the packet consisting of $A(h)$ copies of $\psi h$. We compare packets for the same $Q,\psi,H$ and require every factor that occurs to be primitive of conductor $Q$. This keeps the local factors at primes dividing $Q$ identical. Write $m_A=\sum_h A(h)$. With $\zeta=e^{2\pi i/E}$, define for each $a\in G$
\[
 f_{A,a}(u)=\sum_{h:\,h(a)=\zeta^u} A(h),\qquad
 F_A(a,k)=\sum_{u\bmod E} f_{A,a}(u)\zeta^{ku}.
\]
A rotation of a histogram $f$ is $u\mapsto f(u-s)$ on $\Z/E\Z$. It multiplies its $k$th Fourier coefficient by $\zeta^{ks}$.

\paragraph{The analytic input.}
For a fixed primitive character $\chi$ and fixed $x>1$, the Landau explicit formula gives
\[
 \sum_{0<\Im\rho\le T}x^\rho
 =-\frac{T}{2\pi}\,\Lambda(x)\chi(x)+O_{\chi,x}(\log(2T)).
\]
Here the product on the right is defined to be zero unless $x$ is an integral prime power. This follows, including for nonintegral $x$, from Ford--Zaharescu \cite[Lemma 2.1]{FZ}; Banks gives a Dirichlet-character version \cite[Theorem 1.2]{Banks}. Dirichlet GRH gives $\rho=1/2+i\gamma$. Dividing by $T\sqrt{x}$, summing over factors and conjugating at negative frequencies proves that $S_A(t;T)=L_A(t)+O_{A,t}(\log(2T)/T)$, where
\begin{equation}
 L_A(k\log p)=-\frac{\log p}{2\pi p^{|k|/2}}\,\psi(p)^k F_A(p,k)
 \quad(p\nmid Q,\ k\in\Z\setminus\{0\}). \label{eq:landau}
\end{equation}
All other nonzero frequencies have limit zero. We write $p$ for its residue in $G$ when evaluating a character. The qualitative fixed-frequency limit also follows from Ford--Soundararajan--Zaharescu \cite[Theorem 5 and Corollary 1]{FSZ} by taking a single Fourier mode. The critical-line result supplies the direct error estimate and the growing-bandwidth assertion below.

\begin{theorem}[Local classification]\label{th:local}
The following conditions are equivalent.
\begin{enumerate}
\setlength\itemsep{0pt}
\item The limits of \eqref{eq:observable} agree for $A,B$ for all balanced tuples.
\item They agree for tuples at a single prime, of orders $2\le d\le\max(2,E)$, using one prime in each reduced residue class modulo $Q$.
\item For each $a\in G$, $f_{B,a}$ is a rotation of $f_{A,a}$.
\end{enumerate}
For the tests in condition 2 the frequencies are $k_j\log p$, with nonzero integer $k_j$ of sum zero. Finitely many choices suffice.
\end{theorem}

\newpage
\section{Proof of the classification and global recovery}
\begin{proof}[Proof of Theorem~\ref{th:local}]
Suppose condition 3 holds. If any slot of a tuple is outside the support allowed in \eqref{eq:landau}, the product has zero limit for both packets. Otherwise write $t_j=k_j\log p_j$. Unique factorization applied to $\prod_j p_j^{k_j}=1$ gives $\sum_{p_j=p} k_j=0$ separately at every prime. The rotation factors for each prime therefore multiply to one. This proves condition 1, which gives condition 2 by restriction. Dirichlet's theorem provides the requisite prime representatives.

For the converse, the pair $(E\log p,-E\log p)$ determines $m_A^2$, hence $m_A$. Fix a class $a$ and its chosen prime $p$, and abbreviate $f=f_{A,a}$, $F(k)=F_A(a,k)$. After division by the known factors in \eqref{eq:landau}, the tests give every product $\prod_{j=1}^d F(b_j)$ with $\sum_j b_j=0\bmod E$ and $2\le d\le E$. Indeed, lift the residues $b_j$ to integers of sum zero. Adding multiples of $E$ to the first $d-1$ entries and subtracting their sum from the last lets one avoid zero in every slot. Choose one such lift for each residue tuple. Degree one requires only $F(0)=m_A$.

For any real vector $y$ indexed by $\Z/E\Z$, form $z_s=\sum_u y_u f(u-s)$. Fourier inversion and orthogonality show that $\sum_s z_s^d$ is a linear combination of the products just recovered: summation over $s$ kills every term of nonzero total frequency modulo $E$. Thus the first $E$ power sums of the $E$ numbers $z_s$ agree for $A$ and $B$. Newton's identities give equality of their multisets. Choose $y_u=(m_A+1)^u$ using representatives $0\le u<E$. Each $z_s$ now encodes one rotation as base-$(m_A+1)$ digits. Equality of these integers gives condition 3. This proof also covers $E=1$.
\end{proof}

\begin{corollary}[Coprime packet size]\label{th:glue}
If $\gcd(m_A,E)=1$, equality of all balanced data is equivalent to $B$ being obtained by multiplying every member of $A$ by a single character from $H$.
\end{corollary}
\begin{proof}
For the twists alone put $D_A=\prod_{h\in H}h^{A(h)}\in H$. The local rotations $v_a\in\muE$ obey $v_a^{m_A}=D_B(a)/D_A(a)$. If $m_Au\equiv1\bmod E$, then $v_a=((D_B/D_A)^u)(a)$ for every $a$. The local sums $F_B(a,1)$ agree everywhere with the sums for this global twist of $A$. Character orthogonality on $G$ gives equality of the multiplicities. Conversely a common twist rotates every local histogram.
\end{proof}

For $E=2$, put $c_A(h)=\sum_{u\in H} A(u)A(hu)$, the multiplicity of $h$ in the multiset of pairwise quotients. Then $\widehat c_A(v)=|\widehat A(v)|^2$ on $\widehat H$. Every $v\in\widehat H$ is evaluation at some $a\in G$: the evaluation image has trivial annihilator in $H$, hence is all of $\widehat H$ by finite duality. A two-bin histogram of known mass is determined up to a swap by the absolute difference of its bins. Consequently
\[
 \boxed{\ \text{all balanced limits agree}\ \Longleftrightarrow\
       \text{all pair limits agree}\ \Longleftrightarrow\ c_A=c_B.\ }
\]
Equality $c_A=c_B$ is called \emph{homometry}. It includes repeated factors.

\paragraph{A smoothed pair-difference consequence.}
There is also a test on bounded gaps with a different normalization. Let $w\in C_c^\infty(\R\setminus\{0\})$ be real with $\int w=0$, and set $P_A(w;T)=\sum_{\gamma,\gamma'\in Z_A(T)}w(\gamma-\gamma')$. For any packets satisfying condition 3 of Theorem~\ref{th:local},
\[
 P_A(w;T)-P_B(w;T)=O_{A,B,w}\bigl(T/(\log T)^{1/3}\bigr).
\]
Indeed, sum the cross-character formula of Ford--Zaharescu \cite[Theorem 2]{FZ} over ordered factors. Its main term depends on the Dirichlet series with coefficients $\Lambda(n)^2|\sum_{\chi\in A}\chi(n)|^2$, where the sum includes repetitions. These coefficients agree at every prime power by local rotation and vanish elsewhere. The continued kernels therefore agree, and summing the finitely many errors gives the assertion. The test $w$ belongs to their class $\mathcal H_0$ by compact support, smoothness and zero integral.


\newpage
\section{Six characters and 28 indistinguishable twist classes}\label{sec:six}
Take $(q_1,q_2,q_3,q_4)=(5,13,17,29)$, so $Q=32045$. At each prime choose a character $\psi_i$ of order four. Set $\psi=\prod_i\psi_i$ by the Chinese remainder theorem and $h_x=\prod_i\psi_i^{2x_i}$ for $x\in V=\F_2^4$. Each $\psi h_x$ is primitive of conductor $Q$, because its local exponents are $1$ or $3$. Consider the two sets of indices where, respectively, the following forms equal one:
\[
 q_A(x)=x_1x_2+x_3x_4,\qquad q_B(x)=x_1x_2+x_3x_4+x_1x_3 .
\]
\begin{theorem}\label{th:six}
These six-factor packets have the same balanced limits at every order. Exactly 448 six-element subsets of $V$ share their data, forming 28 classes under a common Dirichlet twist. For the displayed pair the total multiplicity discrepancy up to height $T$ satisfies
\begin{equation}
 \sum_{0<\gamma\le T}|m_A(\gamma)-m_B(\gamma)|
 \ \ge\ \frac{\sqrt2\log 2}{\pi}T-O(\log(2T)). \label{eq:gap}
\end{equation}
Here $m_A(\gamma)$ is the pooled zero multiplicity at that exact ordinate.
\end{theorem}
\begin{proof}
The Walsh sum of $(-1)^{q_A(x)}$ against $(-1)^{v\cdot x}$ is $4(-1)^{v_1v_2+v_3v_4}$: sum each hyperbolic pair separately. Substitution $x_2'=x_2+x_3$ gives the same formula for $q_B$ with the additional sign $(-1)^{v_2v_4}$. Since the indicator is $(1-(-1)^q)/2$, both sets have size six, and their nontrivial Fourier coefficients all have magnitude two. The quadratic classification applies. Polarization shows the two forms have different bilinear parts, which a translation preserves. A twist relating two nonempty subsets of $\psi H$ has to be the quotient of two members, so translations in $V$ cover all common twists.

Fourier inversion of the squared magnitudes $36,4,\ldots,4$ gives difference counts $6$ at zero and $2$ elsewhere. Equivalently, the 15 unordered pairs of a six-set have distinct nonzero sums. Such a set has no affine dependence of size two or four. Six points in dimension four are affinely dependent, and in characteristic two the support of an affine dependence has even size. Hence all six sum to zero, every five form an affine basis, and the sixth is forced by the first five. Conversely every ordered affine basis extends in this way; the resulting six points have no four-term dependence by independence of the first five. Dividing the number of ordered affine bases by $6!$ gives
$16(15)(14)(12)(8)/6!=448$. A nonzero stabilizing translation would give difference count six at that nonzero shift. Every translation orbit thus has size 16. Allowing repetition adds no further packets in this class, since $c_A(0)=6=m_A$ forces all multiplicities to be zero or one.

For $p\nmid Q$ the sign vector is $(-1)^{v_i}=(p/q_i)$. The displayed Fourier coefficients differ precisely if $v_2=v_4=1$, which holds for one quarter of reduced classes and for $p=2$. Their absolute difference there is four. Apply \eqref{eq:landau} at $\log 2$ and use the triangle inequality before normalization to get \eqref{eq:gap}.
\end{proof}
The same argument on $(T,T+H]$, by subtracting the two endpoint formulas for the complex Fourier sum, gives discrepancy at least $(\sqrt2\log2/\pi)H-O(\log(2T+2H))$. In particular it is positive in every window of length $C\log T$ for a sufficiently large fixed $C$ and large $T$.

\newpage
\section{Uniformity for frequencies chosen at height \texorpdfstring{$T$}{T}}\label{sec:uniform}
Fixed-frequency agreement allows a possible way to distinguish two packets: choose the test after seeing the height, and tune it very close to a prime-power frequency. The next estimate controls all these choices simultaneously.
\begin{theorem}\label{th:uniform}
Suppose the two local histograms are rotations of one another at every reduced residue class, as in Theorem~\ref{th:local}. For fixed $d\ge2$, $a>0$, $0<\delta<2$ and $\eps>0$, as $T\to\infty$,
\[
 \sup_{\substack{\sum_j t_j=0\\ a\le |t_j|\le (2-\delta)\log T}}
 |\cB_A(\mathbf t;T)-\cB_B(\mathbf t;T)|
   \ll_{A,B,d,a,\delta,\eps} T^{-\delta/2+\eps}.
\]
The bandwidth exponent is independent of the fixed order $d$.
\end{theorem}
\begin{proof}
Put $X=T^{2-\delta}$, $x=e^{|t|}$, $s=\operatorname{sgn}(t)$, and let $n$ be a nearest integer to $x$, resolving ties consistently. Write $K_T(u)=T^{-1}\int_0^T e^{ivu}\,dv$. The uniform Landau formula \cite[Lemma 2.1]{FZ} and GRH approximate each $S_A(t;T)$, for $e^a\le x\le X$, by zero if $n$ is not an unramified prime power, and otherwise, for $n=p^k$, by
\begin{equation}
 -\frac{\log p}{2\pi\sqrt x}\,\psi(p)^{sk} F_A(p,sk)\, K_T(t-s\log n). \label{eq:sinc}
\end{equation}
The absolute error in either case is
$O_{A,a,\eta}(X^{1/2+\eta}\log(2XT)/T)$ for every $\eta>0$. Indeed, the uniform formula approximates $\sum x^\rho$ by the nearest-integer term and has error
$O_{\chi,\eta}(x^{1+\eta}\log(2x)+x^{1+\eta}\log T+\log T/\log x)$ for a Dirichlet factor. Divide by $T\sqrt x$. Complex conjugation gives the displayed signed kernel. Each leading term is $O_A(1)$ because $|K_T|\le1$ and $\log n/\sqrt{x}$ is bounded. Choose $\eta$ small enough that the error tends to zero. A telescoping product propagates the same error through $d$ factors.

If every nearest integer is an unramified prime power, set $U=\prod_{t_j>0} n_j$ and $V=\prod_{t_j<0} n_j$. When $U=V$, the signed exponents at each prime add to zero, so all local rotation factors cancel. The products of the leading terms are identical, including the kernels. When $U\ne V$, balance of the $t_j$ gives at least one detuning $|t_j-s_j\log n_j|\ge|\log(U/V)|/d$. Use $|K_T(u)|\le\min(1,2/(T|u|))$, and $n_j\asymp x_j$, to bound either leading product by
\[
 \ll_{A,B,d}\frac{(\log(2X))^d}{T\sqrt{UV}\,|\log(U/V)|}
 \ll_{A,B,d}\frac{(\log(2X))^d}{T}.
\]
For the last inequality, distinct positive integers satisfy
$|\log(U/V)|\ge1/(1+\min(U,V))$ and $\sqrt{UV}\ge\min(U,V)$. This cancellation controls near coincidences between large products independently of their number of prime factors. If any slot has zero main term, both products already vanish. Substitution of $X$ and a sufficiently small $\eta$ proves the claim.
\end{proof}
Ordered tuples of distinct zero occurrences give the same limits. Uniformly in the frequencies, terms with a repeated occurrence number $O_d((T\log T)^{d-1})$, hence contribute $O_d(T^{-1}\log^{d-1}T)$ for fixed packets. The estimate in Theorem~\ref{th:uniform} therefore applies in that convention as well.

\newpage
\section{One complex cubic at odd prime exponent}\label{sec:cubic}
When the twist exponent is an odd prime $\ell$, a single test at each chosen prime records substantially more than the pair measurements. The statement below assumes known packet size $m$; it can always be obtained by the pair test in the proof of Theorem~\ref{th:local}.
\begin{theorem}\label{th:cubic}
Let $E=\ell$ be an odd prime. For each $a\in G$, choose one prime $p\equiv a\pmod Q$. The limiting measurement at $(\log p,\log p,-2\log p)$ determines $f_{A,a}$ up to cyclic rotation. If $\ell\nmid m$, these measurements determine the entire packet up to a common twist.
\end{theorem}
\begin{proof}
For $f=f_{A,a}$ write $\alpha_f=\sum_{u=0}^{\ell-1} f(u)\zeta^u\in\Z[\zeta]$. Let $\sigma(\zeta)=\zeta^2$. Removing the known nonzero real factor $-(\log p)^3/((2\pi)^3p^2)$ from the cubic gives
\[
 C(f)=\alpha_f^2\overline{\sigma(\alpha_f)} .
\]
The factors from $\psi$ cancel. If $C(f)=0$, the minimal polynomial $1+X+\cdots+X^{\ell-1}$ shows that all entries of $f$ are equal, and $m$ determines their value. Otherwise compare $g$ with $C(g)=C(f)$ and let $b=\alpha_g/\alpha_f$. Then $\sigma(b)=\overline b^{-2}$. Complex conjugation commutes with $\sigma$. Along its finite orbit, the radii $r_j=|\sigma^j(b)|$ obey $r_{j+1}=r_j^{-2}$. If the orbit closes after $h=\ord_\ell(2)$ steps, $r_0=r_0^{(-2)^h}$, whence $r_0=1$. Thus $\sigma(b)=b^2$ and $b^{2^h-1}=1$.

Every root of unity of odd order $n$ in $\Q(\zeta)$ has $n\mid\ell$. Indeed adjoining a primitive $n$th root gives the cyclotomic field of order $\operatorname{lcm}(n,\ell)$, and its degree exceeds $\ell-1$ if this odd least common multiple exceeds $\ell$. Hence $b=\zeta^s$. Subtract the corresponding rotation of $f$ from $g$. Its Fourier coefficient at one vanishes, so its entries are constant by the same minimal polynomial argument. Equal mass forces that constant to be zero. The last claim follows from Corollary~\ref{th:glue}.
\end{proof}

\paragraph{Precision of the cubic.}
A nonzero difference $D=C(f)-C(g)$ is an algebraic integer. All its conjugates have modulus at most $2m^3$. Pairing the identity embedding with complex conjugation in its nonzero integer norm gives
\[
 |D|\ge (2m^3)^{-(\ell-3)/2}.
\]
Thus absolute error below half this separation suffices for finite enumeration of the local rotation class. This is a bound in the normalized cubic variable. A numerical certificate from truncated zero lists additionally needs an error enclosure for that truncation.

\paragraph{Why a complex measurement matters.}
The local histogram with one copy at $0$ and two at $1$, and its reflection with two at $-1$, have identical quadratic data at every frequency. They are distinct up to rotation for every odd $\ell$, since their unique bin of size one is at zero. The cubic distinguishes them. Jaming--Kolountzakis \cite{JK} treat reconstruction from the full triple correlation. Casper--Orozco \cite[Theorem 3.1]{CO} use doubling orbits with both quadratic and cubic data. Here integrality and the finite orbit of the radii allow one complex cubic to suffice.

\newpage
\section{Larger families and the source of the ambiguity}
One can produce many more collisions, including pairs that remain inequivalent under affine changes of the character labels. Take $2r$ distinct primes $q_i\equiv1\pmod4$ and the same quartic-character construction as in Section~\ref{sec:six}. Identify the quadratic twists with $\F_2^r\times\F_2^r$. For any Boolean function $g:\F_2^r\to\F_2$ with $g(0)=0$, take the support
\[
 A_g=\{(x,y):\ x\cdot y+g(y)=1\}.
\]
These are instances of the classical Maiorana--McFarland construction of bent functions, related to Hadamard difference sets \cite{Dillon,Carlet}. Its Walsh sum has a short proof:
\[
 \sum_{x,y}(-1)^{x\cdot y+g(y)+u\cdot x+v\cdot y}
   =2^r(-1)^{g(u)+v\cdot u}.
\]
The inner sum over $x$ vanishes unless $y=u$. Every packet therefore has size $2^{2r-1}-2^{r-1}$ and nontrivial indicator Fourier coefficients of magnitude $2^{r-1}$. All have identical balanced data, with the uniform bound of Theorem~\ref{th:uniform}.

A translate of the support of $x\cdot y+g(y)$ by $(a,b)$ is the support of $x\cdot y+x\cdot b+a\cdot y+a\cdot b+g(y+b)$. Comparing the coefficients of $x$ forces $b=0$ if this is $A_h$. Then $h(y)=g(y)+a\cdot y$. Exactly $2^r$ choices lie in each common-twist class of this family. The resulting number of classes is
\[
 2^{\,2^r-1-r}.
\]
For $r\ge3$, take $g=0$ and $h(y)=y_1y_2y_3$. The Boolean polynomials defining the two supports have different algebraic degrees. Invertible affine substitution preserves degree, by applying the degree bound to the substitution and its inverse. This pair remains distinct even if the labels can be changed affinely.

\paragraph{Relation to the requested sources.}
The shortened Lean development proves the open-strip critical-line assertion for all Dirichlet characters \cite{Lean}, with independent verification for the compact edition \cite{Rub}. Its one-sided prime-power reduction uses a quadratic twist to control competing zero contributions. ProveIt's collection \cite{PI} works out many transfers; one packet recovers a modulus and a twist prime from a signed zero list using absolute-height probes. Our question changes the observation to a function of differences of pooled heights. Formula~\eqref{eq:landau} explains exactly where the resulting loss of information occurs. The linked overview \cite{Wiki} describes the prime--zero explicit-formula connection. The uniform Landau estimate, finite cyclic correlation theory and bent-function construction are credited ingredients. The classification for Dirichlet zero packets, its uniform frequency range, and its explicit ambiguity and recovery examples are the results derived here. We found no claim of this combination in the cited sources. Historical priority remains open.

\paragraph{Further questions.}
Can Theorem~\ref{th:uniform} be extended to the scale $e^{|t|}\asymp T^2$? What is the smallest packet size yielding a nontrivial ambiguity in each exponent? For a general composite exponent, which explicit list of cubic and higher measurements is shortest at a single unramified prime?

\newpage
\begin{thebibliography}{99}
\raggedright
\bibitem{Lean} L. Sernau and collaborators, \emph{GRH for Dirichlet $L$-functions, in Lean 4}, commit \texttt{72fead102f663d5c7a407964a13dc0cf9d977207}. \href{https://github.com/lsernau/grh-lean/blob/72fead102f663d5c7a407964a13dc0cf9d977207/grh/src/Statements.lean}{Formal statements}, \href{https://github.com/lsernau/grh-lean/blob/72fead102f663d5c7a407964a13dc0cf9d977207/grh/src/New/Reduction/Zeros/Criterion.lean}{one-sided criterion}, and \href{https://github.com/lsernau/grh-lean}{repository}.
\bibitem{Rub} M. Rubinstein, \href{https://www.math.uwaterloo.ca/~mrubinst/alpoge_claude_grh/}{\emph{Alp\"oge/Claude proof of the Generalized Riemann Hypothesis}}, independent verification reports, including the shortened edition, 26 November 2026.
\bibitem{PI} V. Reshetnikov, \href{https://github.com/VladimirReshetnikov/ProveIt/tree/main/Analysis/GeneralizedRiemannHypothesis}{\emph{ProveIt, Generalized Riemann Hypothesis collection}}. In particular, \href{https://github.com/VladimirReshetnikov/ProveIt/blob/main/Analysis/GeneralizedRiemannHypothesis/Research/grh-torsion-and-zero-data/recovery.tex}{\emph{Recovering both arithmetic parameters from a finite signed zero list}}, blob \texttt{de46db208bfb26a1477570c1517f228f6e09ff44}. The collection README describes the editorial status of its 61 packets.
\bibitem{Wiki} Wikipedia contributors, \href{https://en.wikipedia.org/wiki/Riemann_hypothesis}{\emph{Riemann hypothesis}}, sections on the formal proof and the distribution of primes, accessed 26 November 2026.
\bibitem{FZ} K. Ford and A. Zaharescu, \href{https://arxiv.org/abs/1305.2520}{Unnormalized differences between zeros of $L$-functions}, \emph{Compos. Math.} \textbf{151} (2015), 230--252. The uniform estimate is Lemma 2.1; the integrated pair formula is Theorem 2.
\bibitem{Banks} W. D. Banks, \href{https://arxiv.org/abs/2302.07073}{Comparing zeros of distinct Dirichlet $L$-functions}, \emph{Publ. Math. Debrecen} \textbf{106} (2025), 327--335.
\bibitem{FSZ} K. Ford, K. Soundararajan and A. Zaharescu, \href{https://arxiv.org/abs/0805.2745}{On the distribution of imaginary parts of zeros of the Riemann zeta function, II}, \emph{Math. Ann.} \textbf{343} (2009), 487--505.
\bibitem{JK} P. Jaming and M. N. Kolountzakis, \href{https://nyjm.albany.edu/j/2003/9-12.pdf}{Reconstruction of functions from their triple correlations}, \emph{New York J. Math.} \textbf{9} (2003), 149--164.
\bibitem{CO} W. R. Casper and B. Orozco, \href{https://arxiv.org/abs/2503.21022}{Reconstructing rational functions on finite abelian groups with higher autocorrelations}, arXiv:2503.21022.
\bibitem{Dillon} J. F. Dillon, \href{https://drum.lib.umd.edu/items/d8150b1f-1bdd-4c3d-9987-c14a95251edc}{\emph{Elementary Hadamard Difference Sets}}, Ph.D. thesis, University of Maryland, 1974.
\bibitem{Carlet} C. Carlet, \href{https://www.math.univ-paris13.fr/~carlet/book-fcts-Bool-vect-crypt-codes.pdf}{\emph{Boolean Functions for Cryptography and Coding Theory}}, Cambridge University Press, 2021.
\end{thebibliography}

\paragraph{Reproducibility.}
The first Python program enumerates all 8,008 six-subsets and checks all 21,504 reduced classes for the conductor-$32045$ example. The second uses exact cyclotomic arithmetic on all 1,100 histograms of sizes zero through five for $\ell=3,5,7$. Their recorded outputs accompany the source. The general statements and the estimates for zeros are proved in the text. They have not been formalized in Lean.
\end{document}
